\section{Limitations}

The current evidence supports a scoped practical-use claim, not a broad claim
of synthetic-data superiority. DataFrameSampler should be evaluated as a simple
and inspectable generator for examples, tests, demonstrations, and exploratory
workflows. It should not be presented as a clinical simulator, a regulated
statistical disclosure-control method, or a replacement for validated
domain-specific synthetic-data systems.

The anonymization layer does not provide differential privacy, formal
de-identification, $k$-anonymity, HIPAA compliance, GDPR compliance, or
protection against linkage attacks. It checks exact and normalized reuse of
selected source values after replacement and generation. Columns that are not
selected for replacement may still emit observed source values because the
decoder samples from values observed in the fitted data. Because sampling is
based on nearest-neighbour chains, very small datasets or highly distinctive
rows may also make generated rows close to particular source records. This
memorization risk must be considered separately from exact identifier
replacement.

The LLM-assisted configuration mode is a convenience interface. It can reduce
the burden of choosing vectorization columns, sampled columns, embedding
methods, and nearest-neighbour backends, but its recommendations remain
heuristic. Opaque column names, coded values, missing domain context, or poor
examples can produce weak configurations. User review is therefore part of the
intended workflow.

Finally, the method's behaviour depends on the quality of the latent
representation and neighbour graph. High-cardinality categorical columns, sparse
bins, rare categories, weak one-dimensional categorical embeddings, or approximate-neighbour errors may
produce implausible rows. These cases should be treated as boundary conditions
for the final claim rather than hidden as implementation details.

\begin{table*}
\caption{Limitations and allowable conclusion scope under the current evidence. Takeaway: the paper's strongest defensible conclusion is practical, inspectable example generation under explicit boundaries.}
\label{tab:limitations-scope}
\begin{adjustbox}{max width=\textwidth}
\begin{tabular}{llll}
\toprule
Area & Known & Unknown & Allowable scope \\
\midrule
Similarity & Public and controlled metrics & Stability on broader task collections & Regime-specific mixed-type evidence \\
Utility & Held-out real utility tests & General predictive advantage & Utility only where lift is observed \\
Privacy & Nearest-neighbour and overlap diagnostics & Linkage risk under a threat model & No formal privacy claim \\
Usability & Notebook runtime and configs & User study or setup-time distribution & Practical workflow claim only \\
\bottomrule
\end{tabular}
\end{adjustbox}
\end{table*}

